\documentclass[UTF8,fontset=none,11pt,a4paper]{ctexart}
\usepackage{ipara-handbook}
\handbooksetup{NET-08 应用层服务语义}{408 · 计算机网络 NET-08}{DNS · HTTP · FTP · Mail}
\providecommand{\tightlist}{\setlength{\itemsep}{0pt}\setlength{\parskip}{0pt}}
\providecommand{\passthrough}[1]{#1}
\begin{document}
\handbooktitle{应用层服务语义}{把字节组织成名字、资源与操作}{408 · NET-08}
\section{0. 本册定位}\label{0-ux672cux518cux5b9aux4f4d}

本 Topic 回答：端到端传输只提供 byte stream 或
datagram，应用怎样定义对象、名字、操作、响应和状态，使不同程序能够共同完成
DNS、Web、文件与邮件服务？

本册拥有 application architecture、DNS、HTTP semantics、FTP
与电子邮件的服务分工，以及 application protocol 对 transport service
的选择逻辑。

DHCP 虽以 UDP application protocol
形式运行，但在本项目中由\href{../04_IP地址_子网与分组转发/README.md}{IP
配置与转发}拥有，因为其主要输出是主机进入 IP 世界所需的启动配置。

\section{1. 根本问题：transport 不知道``这段 byte
意味着什么''}\label{1-ux6839ux672cux95eeux9898transport-ux4e0dux77e5ux9053ux8fd9ux6bb5-byte-ux610fux5473ux7740ux4ec0ux4e48}

TCP 可以保证 byte 顺序，却不知道这段 byte 是域名查询、HTTP
request、邮件命令还是文件内容。Application protocol 必须定义：

\[
\boxed{
\text{Application object}
\to \text{Name/identifier}
\to \text{Operation}
\to \text{Message representation}
\to \text{Response/state transition}
}
\]

它不仅规定报文格式，还规定动作的语义与交互顺序。相同 transport
可以承载完全不同服务；相同服务语义也可以演进出不同 transport bindings。

\section{2.
应用架构：谁长期持有状态}\label{2-ux5e94ux7528ux67b6ux6784ux8c01ux957fux671fux6301ux6709ux72b6ux6001}

\subsection{2.1 Client/Server}\label{21-clientserver}

server 通常在稳定地址/名字上等待请求并持有共享服务状态；client
主动发起交互。集中服务便于管理、一致性和发现，代价是容量、故障和运营中心化。

\subsection{2.2 Peer-to-Peer}\label{22-peer-to-peer}

peer 同时消费和提供资源，容量可随参与者增加，但发现、可用性、信任、NAT
穿越和一致性更复杂。

C/S 与 P2P 是状态与职责的组织方式，不等同于
TCP/UDP，也不由``是否有服务器进程''一个表象决定。

\section{3.
DNS：把一个集中表拆成层次化授权系统}\label{3-dnsux628aux4e00ux4e2aux96c6ux4e2dux8868ux62c6ux6210ux5c42ux6b21ux5316ux6388ux6743ux7cfbux7edf}

\subsection{\texorpdfstring{3.1 为什么单一 \texttt{hosts}
文件失败}{3.1 为什么单一 hosts 文件失败}}\label{31-ux4e3aux4ec0ux4e48ux5355ux4e00-hosts-ux6587ux4ef6ux5931ux8d25}

全网共享一份名字表会遇到：更新冲突、分发延迟、单点容量与管理边界。DNS
的生成链是：

\[
\boxed{
\text{Global flat file fails}
\to \text{Hierarchical namespace}
\to \text{Delegated zones}
\to \text{Distributed authoritative servers}
\to \text{Resolver + cache}
}
\]

Domain namespace 是名字树；zone 是某个 authority
实际管理的数据边界。树的结构与管理分区相关，但 domain
不等于一台主机或一个物理网络。

\subsection{3.2 四类对象}\label{32-ux56dbux7c7bux5bf9ux8c61}

\begin{itemize}
\tightlist
\item
  stub resolver：应用/OS 发起查询的本地接口；
\item
  recursive resolver：代用户完成后续查询并缓存；
\item
  authoritative name server：对某 zone 的数据负责；
\item
  root/TLD/other delegated servers：通过 referral 把查询逐步引向
  authority。
\end{itemize}

\subsection{3.3 一次 cache miss
的解析轨迹}\label{33-ux4e00ux6b21-cache-miss-ux7684ux89e3ux6790ux8f68ux8ff9}

\begin{lstlisting}
application asks local resolver
-> recursive resolver checks cache
-> root referral identifies TLD servers
-> TLD referral identifies authoritative servers
-> authoritative answer returns records
-> resolver caches answer/referrals under TTL
-> result returns to application
\end{lstlisting}

迭代查询返回``我不知道最终答案，但下一步问谁''；递归查询把继续查找的责任交给被询问方。实际
client-to-recursive 常请求递归，resolver 对 authority 链执行迭代式追踪。

\subsection{3.4 Cache
的正确性边界}\label{34-cache-ux7684ux6b63ux786eux6027ux8fb9ux754c}

cache 用旧信息换低延迟和低查询负载。TTL 限制可复用时间，但 DNS
不是瞬时全局一致系统：更新在 TTL 生命周期内可能与缓存旧值并存。

DNS 并非只能用 UDP。经典普通查询常用 UDP；响应截断、zone transfer
等情形可使用 TCP，现代部署还存在加密传输。408
做题要服从题设，不能把``DNS=UDP''写成协议语义。

\section{4.
HTTP：操作资源的统一语义}\label{4-httpux64cdux4f5cux8d44ux6e90ux7684ux7edfux4e00ux8bedux4e49}

\subsection{4.1 WWW：分布式超文档系统}

WWW 不是某一个 protocol 或“所有 Internet 服务”的同义词。它至少包含：browser/client、Web server、由 URL/URI 标识的 resource、HTML 等 document/representation，以及把文档与其他资源连接起来的 hyperlink。HTTP 定义 client 与 server 操作这些资源的请求/响应语义；DNS 可把 URL 中的 host name 解析为 IP；transport 再承载 HTTP message。

\[
\boxed{\text{URL identifies resource}\to\text{DNS locates host}\to
\text{HTTP requests representation}\to\text{HTML links reveal dependencies}}
\]

“网页”常是一个主文档与图片、样式、脚本等多个对象组成的依赖图，不是一条 HTTP response 必然完整返回的单文件。URL 提供标识与定位语法，HTML 提供超文本表示，HTTP 提供操作/传输语义；三者不能互相替代。

HTTP 的核心对象不是``网页文件''，而是由 URI 标识的
resource。消息传输的是 resource 的 representation，method 表达 client
希望对目标执行的语义。

\[
\boxed{
\text{Target resource}
\to \text{Method semantics}
\to \text{Request metadata/content}
\to \text{Status}
\to \text{Representation/metadata}
}
\]

\subsection{4.2 Resource 不等于
representation}\label{41-resource-ux4e0dux7b49ux4e8e-representation}

同一 resource 可根据时间、语言、编码或内容协商产生不同
representation。response body 是某次选定表示，不等于资源本身的全部身份。

\subsection{4.3 Method
属性改变重试决策}\label{42-method-ux5c5eux6027ux6539ux53d8ux91cdux8bd5ux51b3ux7b56}

\begin{itemize}
\tightlist
\item
  safe：语义上只读，例如 GET/HEAD；
\item
  idempotent：重复执行与执行一次具有相同预期 effect，例如 PUT、DELETE
  以及 safe methods；
\item
  POST 通常不假定 idempotent。
\end{itemize}

这些属性不是礼貌标签，而是 cache、proxy、client
在失败后能否自动重试的重要决策依据。

\subsection{4.4 Stateless
不等于``服务没有状态''}\label{43-stateless-ux4e0dux7b49ux4e8eux670dux52a1ux6ca1ux6709ux72b6ux6001}

HTTP stateless 表示每个 request 的语义不依赖 server 必须记住此前
protocol request context。应用仍可通过 cookie、token、database 和
server-side session 建立业务状态。

\subsection{4.5
语义与连接版本分离}\label{44-ux8bedux4e49ux4e0eux8fdeux63a5ux7248ux672cux5206ux79bb}

\begin{itemize}
\tightlist
\item
  HTTP/1.1 和 HTTP/2 通常在 TCP 上承载；
\item
  HTTP/3 使用 QUIC over UDP；
\item
  核心 method、status、representation 和 cache semantics
  在版本间保持共同定义。
\end{itemize}

因此``HTTP 使用 TCP''是特定版本/部署的 transport mapping，不是 HTTP
resource semantics 本身。

\section{5. Web
性能从依赖图生成}\label{5-web-ux6027ux80fdux4eceux4f9dux8d56ux56feux751fux6210}

取得一个 URL 可能依次触发：DNS、transport/security establishment、HTTP
request/response、HTML parsing、更多 subresource requests。

总时间不能只数``几个 RTT''，必须先画依赖关系：哪些步骤串行、哪些
connection 可复用、哪些 object 可并行、哪些命中 cache。

Persistent connection 减少重复建立连接；multiplexing 允许多个
request/response 共享连接并发推进；cache 用 freshness/validation
状态减少传输。三者解决的瓶颈不同。

\section{6.
FTP：控制状态与数据传输分离}\label{6-ftpux63a7ux5236ux72b6ux6001ux4e0eux6570ux636eux4f20ux8f93ux5206ux79bb}

FTP 的母问题是跨异构系统可靠操作和传输文件。它使用持久 control
connection 交换命令与会话状态，并为目录/文件数据建立独立 data
connection。

这种 out-of-band control
让命令不被大文件数据阻塞，也带来额外连接管理、主动/被动模式和
NAT/firewall 适配成本。FTP 的两条 TCP connections 是应用层设计，不是 TCP
自动生成的``控制通道''。

\section{7.
电子邮件：提交、转发、存储、读取不是一个协议}\label{7-ux7535ux5b50ux90aeux4ef6ux63d0ux4ea4ux8f6cux53d1ux5b58ux50a8ux8bfbux53d6ux4e0dux662fux4e00ux4e2aux534fux8bae}

邮件系统至少有：user agent、mail server、message transfer 和 mailbox
access。

\begin{lstlisting}
sender user agent
-> submission / SMTP transfer
-> one or more mail servers
-> recipient mailbox
-> POP3 or IMAP access
-> recipient user agent
\end{lstlisting}

\begin{itemize}
\tightlist
\item
  SMTP 负责提交/服务器间 push 式传送；
\item
  POP3 提供较简单的 mailbox download/access 模型；
\item
  IMAP 更强调服务器端 mailbox 状态与多设备同步；
\item
  MIME 扩展 message representation，使非 ASCII
  内容和附件可被邮件体系承载。
\end{itemize}

``发邮件''和``收邮件''必须分开建模，不能把所有过程归给 SMTP。

\section{8. 408 应用协议流程：消息、状态与停止条件}

\subsection{DNS：递归责任与迭代 referral 分开}

stub 向 recursive resolver 发送 query（name、type、class、transaction ID 等）；resolver 先查 cache。未命中时，它依次向 root、TLD 和 authoritative server 发迭代查询，referral 通过 NS 及必要的 glue/address 信息指出下一位 authority；最终 answer 或 negative result 按 TTL 缓存并返回 stub。收到 response 时必须核对 query identity/匹配字段；TTL 到期后旧 cache 不再直接回答。UDP 是经典普通查询的常见承载，截断或 zone transfer 等分支可转 TCP，不能把承载写成唯一语义。

\subsection{HTTP：请求语义与连接时延分开}

一次交互至少推进：
\[
\boxed{\text{method + target + fields + optional content}
\to\text{server processing}
\to\text{status + fields + optional representation}}.
\]
性能题先画依赖：DNS 是否命中、transport 是否已有连接、连接能否复用、对象是否串行/并行、内容大小和链路速率。Persistent connection 减少重复建连，cache 减少请求/传输，multiplexing 改变并发关系；三者不能用同一个“少一个 RTT”口诀替代。

\subsection{FTP：控制连接决定数据连接怎样建立}

client 先建立持久 control connection 并完成认证/命令。需要目录或文件数据时，主动模式由 server 从约定 data port 连接 client 提供的 endpoint；被动模式由 server 提供监听 endpoint，client 主动连接。一次 data transfer 完成后数据连接可关闭，control state 继续存在。端口常数和命令细节以题设为准，稳定边界是 control/data 两条状态轨迹。

\subsection{Mail：提交、relay 与 access 三段}

sender user agent 把 message 提交给 SMTP server；服务器通过 DNS MX 等信息选择下一跳并用 SMTP relay 推送到 recipient mail server；message 进入 mailbox 后，recipient 再用 POP3/IMAP 等访问协议读取。SMTP reply 驱动提交/转发状态，MIME 只定义可承载的内容表示；它们都不拥有 mailbox access。停止条件应分别写“已被下一 SMTP hop 接受”和“用户已从 mailbox 访问”，不能合并成一次端到端 ACK。

\begin{warnbox}{Owner 边界}
DHCP 虽常列在应用层协议表中，但本项目由 NET04 Own其配置生命周期；NET-I01 只调用 DNS、TCP 和 HTTP 的输出，不重复本节报文机制。
\end{warnbox}

\section{9. 怎样选择 transport
service}\label{8-ux600eux6837ux9009ux62e9-transport-service}

应用选择的不是协议名偏好，而是约束组合：

\begin{longtable}[]{@{}ll@{}}
\toprule\noalign{}
应用需求 & 需要观察的 transport 能力 \\
\midrule\noalign{}
\endhead
\bottomrule\noalign{}
\endlastfoot
不能丢 byte、允许额外延迟 & reliable ordered stream \\
保留 message boundary、自定义恢复 & datagram \\
极低启动延迟 & connection setup cost \\
多路流并发且减少队头阻塞 & multiplexing semantics \\
广播 bootstrap & local datagram/broadcast support \\
安全身份与机密性 & security layer/secure transport \\
\end{longtable}

``实时应用一定用 UDP''\,``可靠应用一定用 TCP''都过度简化。应用可在 UDP
之上实现可靠/拥塞控制，或为穿越部署环境选择
TCP；判断必须回到服务需求与实现责任。

\section{9. 概念边界}\label{9-ux6982ux5ff5ux8fb9ux754c}

\begin{longtable}[]{@{}lclll@{}}
\toprule\noalign{}
概念 A & ≠ & 概念 B & 真正区别与题目信号 & 混淆后果 \\
\midrule\noalign{}
\endhead
\bottomrule\noalign{}
\endlastfoot
Domain name & ≠ & IP address & 人/组织命名与网络定位；DNS 维护映射 & 把
DNS 当逐跳路由 \\
Domain & ≠ & Zone & 名字树节点/子树与 authority 管理边界 &
授权与层次结构混乱 \\
Recursive query & ≠ & Iterative query & 被询问方完成解析；被询问方返回
referral & DNS 报文轨迹画错 \\
Resource & ≠ & Representation & 被标识对象与某次可传输表示 &
缓存和内容协商理解错误 \\
HTTP stateless & ≠ & Application has no state &
协议请求独立与业务状态存在可并存 & 认为登录/cookie 违反 HTTP \\
HTTP semantics & ≠ & HTTP transport version & method/status/resource
与具体连接承载分离 & 把 HTTP/3 误判为非 HTTP \\
SMTP & ≠ & POP3/IMAP & 邮件传送与 mailbox access &
邮件全流程协议归属错误 \\
FTP control connection & ≠ & FTP data connection &
命令状态与内容传输分离 & 端口/连接数量推演错误 \\
\end{longtable}

\section{10.
做题调用协议}\label{10-ux505aux9898ux8c03ux7528ux534fux8bae}

\begin{enumerate}
\def\labelenumi{\arabic{enumi}.}
\tightlist
\item
  写应用对象：name、resource、file、message 还是 lease/config；
\item
  写谁持有 authority/state，谁只是 cache/client；
\item
  画 request/response 或 push/pull 的事件方向；
\item
  DNS 题区分 recursion、iteration、authority、cache TTL；
\item
  HTTP 题区分 resource、representation、method 和 connection；
\item
  邮件题拆成 submission/transfer/storage/access；
\item
  性能题画依赖 DAG，再数不可并行的 RTT/transfer；
\item
  最后检查 transport 选择是否来自应用约束而非口诀。
\end{enumerate}

\section{11. 贯穿母例：输入 URL 后第一个 HTTP request
为什么还不能立刻发}\label{11-ux8d2fux7a7fux6bcdux4f8bux8f93ux5165-url-ux540eux7b2cux4e00ux4e2a-http-request-ux4e3aux4ec0ux4e48ux8fd8ux4e0dux80fdux7acbux523bux53d1}

\begin{lstlisting}
URL identifies scheme + authority + target
-> resolver needs server IP (unless cached)
-> IP forwarding needs next hop and local MAC (unless cached)
-> transport/security context may need establishment
-> HTTP request names resource and method
-> response carries status + selected representation
-> representation may reveal more resource dependencies
\end{lstlisting}

URL、domain、IP、port、MAC
分别服务不同作用域。把它们压成一句``浏览器访问服务器''，会隐藏几乎所有可诊断的状态交接。

\section{12. 高频 First
Divergence}\label{12-ux9ad8ux9891-first-divergence}

\begin{itemize}
\tightlist
\item
  把 DNS 说成从 root 一次返回最终 IP：漏掉 delegation/referral；
\item
  认为 cache 命中永远是最新真相：忽略 TTL 与弱一致时间窗；
\item
  写``DNS 用 UDP，HTTP 用 TCP''作为完整机制：把常见承载当唯一语义；
\item
  把 HTTP stateless 理解成网站不能登录：混淆 protocol context 与
  application state；
\item
  把 response body 等同 resource：混淆对象与 representation；
\item
  用 SMTP 从 mailbox 下载邮件：没有分开 transfer 与 access。
\end{itemize}

\section{13.
一页压缩与复原问题}\label{13-ux4e00ux9875ux538bux7f29ux4e0eux590dux539fux95eeux9898}

\[
\boxed{
\text{Name an application object}
\to \text{Choose operation semantics}
\to \text{Represent in messages}
\to \text{Map onto transport}
\to \text{Evolve application state}
}
\]

\begin{enumerate}
\def\labelenumi{\arabic{enumi}.}
\tightlist
\item
  DNS 怎样用 delegation 替代全球单表？
\item
  cache 为什么同时提高性能并引入旧信息窗口？
\item
  HTTP resource 与 representation 为什么必须分开？
\item
  stateless protocol 怎样支撑有状态登录业务？
\item
  SMTP、POP3/IMAP 和 MIME 分别拥有邮件过程的哪一段？
\end{enumerate}

\section{14.
来源与校正说明}\label{14-ux6765ux6e90ux4e0eux6821ux6b63ux8bf4ux660e}

\begin{itemize}
\tightlist
\item
  归档笔记《应用层-FTP,SMTP,POP3,DHCP》提供传统应用协议覆盖；《应用层-DNS》《应用层-WWW与HTTP》为空文件，未被当作已有正文；
\item
  DNS 的 hierarchy、delegation、resolver、authority 与 cache 边界依据
  \href{https://www.rfc-editor.org/rfc/rfc1034.html}{RFC 1034} 校正；
\item
  HTTP 的 stateless semantics、resource/representation
  和跨版本语义边界依据
  \href{https://www.rfc-editor.org/rfc/rfc9110.html}{RFC 9110} 校正；
\item
  DHCP 已按当前 Ownership 移入 IP Topic，不在应用册重复拥有。
\end{itemize}
\end{document}
